% ---
% title: Waiting for a double six
% tags: independence, complement
% ---

\begin{problem}
	A pair of fair dice is thrown $n$ times; the throws are independent.
	(a) What is the probability that a double six (both dice showing $6$)
	appears at least once?
	(b) How large must $n$ be for this probability to be at least $\frac12$?
\end{problem}

\begin{solution}
	\idea ``At least once'' is hard to count directly, but its complement
	``never'' is easy: by independence the probability of ``no double six in
	$n$ throws'' is a product of $n$ equal factors. Then solve an inequality
	with logarithms. (The same idea as in Exercise~\exref{at-least-one-six}.)
	\step{(a)} In one throw of the pair there are $6\cdot6=36$ equally likely
	outcomes, and exactly one of them, $(6,6)$, is a double six. So in one
	throw
	\[ P(\text{double six})=\frac1{36},\qquad
		P(\text{no double six})=1-\frac1{36}=\frac{35}{36}. \]
	Let $E_i$ be the event that throw $i$ is \emph{not} a double six. The event
	``no double six in $n$ throws'' is the intersection $E_1E_2\cdots E_n$.
	The throws are independent, so the probability of this intersection is the
	product of the probabilities:
	\[ P(E_1E_2\cdots E_n)=P(E_1)P(E_2)\cdots P(E_n)
		=\left(\frac{35}{36}\right)^{n}. \]
	``At least one double six'' is the complement of ``no double six'', so
	\[ P(\text{at least one double six})=1-\left(\frac{35}{36}\right)^{n}. \]
	\step{(b)} We need
	\[ 1-\left(\frac{35}{36}\right)^{n}\ge\frac12
		\quad\Longleftrightarrow\quad
		\left(\frac{35}{36}\right)^{n}\le\frac12. \]
	Take the natural logarithm of both sides. The logarithm is increasing, so
	the inequality keeps its direction:
	\[ n\ln\frac{35}{36}\le\ln\frac12=-\ln2. \]
	Now divide by $\ln\frac{35}{36}$. This number is \emph{negative} (since
	$\frac{35}{36}<1$), so dividing by it flips the inequality:
	\[ n\ge\frac{-\ln2}{\ln\frac{35}{36}}=\frac{\ln2}{\ln\frac{36}{35}}
		\approx\frac{0.6931}{0.02817}\approx24.6. \]
	The number of throws is a whole number, so the smallest possible value is
	$n=25$. Check the two neighbors:
	\[ 1-\left(\frac{35}{36}\right)^{24}\approx0.4914<\frac12,
		\qquad
		1-\left(\frac{35}{36}\right)^{25}\approx0.5055\ge\frac12. \]
	So $24$ throws are not enough, $25$ throws are.
	\begin{answer}
		(a) $1-\left(\frac{35}{36}\right)^{n}$ \quad (b) $n\ge25$
		(with $n=24$: $\approx0.4914$; with $n=25$: $\approx0.5055$)
	\end{answer}
	\begin{remark}
		This is a famous 17th-century gambling question. The Chevalier de Méré
		believed that betting on a double six within $24$ throws was
		favorable. The calculation shows that $24$ throws give a probability
		just below $\frac12$, so the bet loses slightly in the long run.
	\end{remark}
	\begin{remark}
		\textbf{Common mistake:} adding the probabilities,
		$n\cdot\frac1{36}$. This would give $\frac{24}{36}=\frac23$ for
		$n=24$ and more than $1$ for $n\ge37$, which is impossible. The events
		``double six in throw $i$'' are not mutually exclusive, so their
		probabilities may not simply be added.
	\end{remark}
\end{solution}
